Podatki sledilnika pogleda vsebujejo časovne in prostorske odnose med meritvami. 
Časovno-prostorska grafovska predstavitev te odnose eksplicitno kodira s povezavami med vozlišči, medtem ko jih klasična tabelarična predstavitev običajno ne.
V diplomski nalogi smo preverili, ali lahko grafovske nevronske mreže te odnose uporabijo pri binarni klasifikaciji valence.
Časovna okna meritev sledilnika pogleda smo predstavili kot časovno-prostorske grafe, v katerih vozlišča ustrezajo meritvam, povezave pa opisujejo časovno bližino, prostorsko bližino in pripadnost isti fiksaciji.
Grafovske modele smo primerjali z izhodiščnima klasifikatorjema, klasičnimi modeli strojnega učenja ter prednaučenima kodirnikoma signalov.
Modele smo ovrednotili s 7-kratnim prečnim preverjanjem po subjektih na štirih množicah signalov: \emph{samo pogled}, \emph{samo zenici}, \emph{pogled in zenici} ter \emph{vsi signali}.
Rezultati kažejo, da so grafovske predstavitve informativne, saj grafovski modeli pri vseh množicah signalov presežejo naključni in večinski klasifikator in so po rezultatih primerljivi z uveljavljenimi modeli.
Najboljši grafovski model doseže $65{,}5~\%$ točnost in $64{,}4~\%$ makro F1 pri množici \emph{pogled in zenici}, vendar ostane pod najboljšim rezultatom modela \ModelName{SVM}, ki doseže $70{,}5~\%$ točnost in $69{,}7~\%$ makro F1.
V izbranih eksperimentih so najuspešnejši klasični modeli z ročno oblikovanimi značilkami, grafovsko modeliranje pa ostaja smiselna raziskovalna smer za podatke sledilnika pogleda.
